\fancyhead[LE,RO]{\textsf{\fontsize{9}{11}\selectfont CHƯƠNG 1\hspace{1.5em}ÔN TẬP}}
\fancyhead[RE,LO]{}
\fancyfoot[C]{}
\fancyfoot[LE,RO]{\textsf{\textbf{\thepage}}}

\noindent
% Banner tiêu đề REVIEW
\begin{tcolorbox}[
    enhanced,
    colback=stewartblue!12,
    colframe=white,
    boxrule=0pt,
    arc=0pt,
    outer arc=0pt,
    left=0pt,
    right=0pt,
    top=0pt,
    bottom=0pt,
    boxsep=0pt,
    height=28pt
]
    \begin{tikzpicture}[remember picture, overlay]
        \node[fill=stewartblue, minimum width=28pt, minimum height=28pt, anchor=west] at (0,0) {
            \color{white}\fontsize{18}{20}\selectfont\textbf{\textsf{1}}
        };
    \end{tikzpicture}%
    \hspace{38pt}{\fontsize{16}{18}\selectfont\textbf{\textsf{\color{black}ÔN TẬP}}}
\end{tcolorbox}

\vspace{0.98em}

% SECTION: KIỂM TRA KHÁI NIỆM
\noindent
{\fontsize{11}{13}\selectfont\textbf{\textsf{\color{stewartcyan}KIỂM TRA KHÁI NIỆM}}}
\hfill
{\fontsize{8.5}{10}\selectfont\textsf{Đáp án phần Kiểm tra khái niệm có tại StewartCalculus.com.}}

\vspace{-0.4em}
\noindent{\color{stewartcyan}\rule{\textwidth}{0.8pt}}
\vspace{0.49em}

\begin{multicols}{2}
\raggedcolumns
\begin{enumerate}[label=\textbf{\arabic*.}, leftmargin=1.5em, itemsep=1.2em, topsep=0pt]
    \item 
    \begin{enumerate}[label=\textsf{(\alph*)}, leftmargin=1.5em, itemsep=0.3em]
        \item Hàm số là gì? Tập xác định và tập giá trị của nó là gì?
        \item Đồ thị của một hàm số là gì?
        \item Làm thế nào để nhận biết một đường cong đã cho có phải là đồ thị của một hàm số hay không?
    \end{enumerate}

    \item Trình bày bốn cách biểu diễn một hàm số. Minh họa lời giải thích của bạn bằng các ví dụ.

    \item 
    \begin{enumerate}[label=\textsf{(\alph*)}, leftmargin=1.5em, itemsep=0.3em]
        \item Hàm số chẵn là gì? Làm thế nào để biết một hàm số là chẵn khi nhìn vào đồ thị của nó? Nêu ba ví dụ về hàm số chẵn.
        \item Hàm số lẻ là gì? Làm thế nào để biết một hàm số là lẻ khi nhìn vào đồ thị của nó? Nêu ba ví dụ về hàm số lẻ.
    \end{enumerate}

    \item Hàm số tăng (đồng biến) là gì?

    \item Mô hình toán học là gì?

    \item Cho một ví dụ về từng loại hàm số sau.
    \begin{enumerate}[label=\textsf{(\alph*)}, leftmargin=1.5em, itemsep=0.2em]
        \item Hàm tuyến tính
        \item Hàm lũy thừa
        \item Hàm mũ
        \item Hàm bậc hai
        \item Đa thức bậc 5
        \item Hàm phân thức hữu tỉ
    \end{enumerate}

    \item Vẽ phác bằng tay, trên cùng một hệ trục tọa độ, đồ thị của các hàm số sau.
    \begin{enumerate}[label=\textsf{(\alph*)}, leftmargin=1.5em, itemsep=0.2em]
        \item $f(x) = x$
        \item $g(x) = x^2$
        \item $h(x) = x^3$
        \item $j(x) = x^4$
    \end{enumerate}

    \item Vẽ phác bằng tay dạng đồ thị của mỗi hàm số sau.
    \begin{multicols}{3}
    \begin{enumerate}[label=\textsf{(\alph*)}, leftmargin=1.5em, itemsep=0.3em]
        \item $y = \sin x$
        \item $y = \tan x$
        \item $y = e^x$
        \item $y = \ln x$
        \item $y = 1/x$
        \item $y = |x|$
        \item $y = \sqrt{x}$
        \item $y = \tan^{-1} x$
    \end{enumerate}
    \end{multicols}

    \item Giả sử hàm số $f$ có tập xác định $A$ và $g$ có tập xác định $B$.
    \begin{enumerate}[label=\textsf{(\alph*)}, leftmargin=1.5em, itemsep=0.3em]
        \item Tập xác định của $f + g$ là gì?
        \item Tập xác định của $fg$ là gì?
        \item Tập xác định của $f/g$ là gì?
    \end{enumerate}

    \item Hàm hợp $f \circ g$ được định nghĩa như thế nào? Tập xác định của nó là gì?

    \item Giả sử đồ thị của $f$ đã biết. Hãy viết phương trình cho mỗi đồ thị nhận được từ đồ thị của $f$ như sau.
    \begin{enumerate}[label=\textsf{(\alph*)}, leftmargin=1.5em, itemsep=0.25em]
        \item Dịch lên trên 2 đơn vị.
        \item Dịch xuống dưới 2 đơn vị.
        \item Dịch sang phải 2 đơn vị.
        \item Dịch sang trái 2 đơn vị.
        \item Lấy đối xứng qua trục $x$.
        \item Lấy đối xứng qua trục $y$.
        \item Dãn theo phương thẳng đứng với hệ số 2.
        \item Co theo phương thẳng đứng với hệ số 2.
        \item Dãn theo phương ngang với hệ số 2.
        \item Co theo phương ngang với hệ số 2.
    \end{enumerate}

    \item 
    \begin{enumerate}[label=\textsf{(\alph*)}, leftmargin=1.5em, itemsep=0.3em]
        \item Hàm số đơn ánh (một-đối-một) là gì? Làm thế nào để biết một hàm số là đơn ánh khi nhìn vào đồ thị của nó?
        \item Nếu $f$ là hàm số đơn ánh, thì hàm ngược $f^{-1}$ của nó được định nghĩa như thế nào? Làm thế nào để suy ra đồ thị của $f^{-1}$ từ đồ thị của $f$?
    \end{enumerate}

    \item 
    \begin{enumerate}[label=\textsf{(\alph*)}, leftmargin=1.5em, itemsep=0.3em]
        \item Hàm sin ngược $f(x) = \sin^{-1} x$ được định nghĩa như thế nào? Tập xác định và tập giá trị của nó là gì?
        \item Hàm cosin ngược $f(x) = \cos^{-1} x$ được định nghĩa như thế nào? Tập xác định và tập giá trị của nó là gì?
        \item Hàm tang ngược $f(x) = \tan^{-1} x$ được định nghĩa như thế nào? Tập xác định và tập giá trị của nó là gì?
    \end{enumerate}
\end{enumerate}
\end{multicols}

\vspace{1.23em}

% SECTION: CÂU HỎI ĐÚNG-SAI
\noindent
{\fontsize{11}{13}\selectfont\textbf{\textsf{\color{stewartcyan}CÂU HỎI ĐÚNG--SAI}}}

\vspace{-0.4em}
\noindent{\color{stewartcyan}\rule{\textwidth}{0.8pt}}
\vspace{0.49em}

\noindent
{\fontsize{9.5}{12}\selectfont
Hãy xác định mỗi khẳng định sau là đúng hay sai. Nếu đúng, hãy giải thích lý do. Nếu sai, hãy giải thích tại sao hoặc đưa ra một ví dụ bác bỏ khẳng định đó.
}

\vspace{0.49em}

\begin{multicols}{2}
\begin{enumerate}[label=\textbf{\arabic*.}, leftmargin=1.5em, itemsep=1.1em, topsep=0pt]
    \item Nếu $f$ là một hàm số, thì $f(s + t) = f(s) + f(t)$.

    \item Nếu $f(s) = f(t)$, thì $s = t$.

    \item Nếu $f$ là một hàm số, thì $f(3x) = 3f(x)$.

    \item Nếu hàm số $f$ có hàm ngược và $f(2) = 3$, thì $f^{-1}(3) = 2$.

    \item Một đường thẳng đứng cắt đồ thị của một hàm số tại nhiều nhất một điểm.

    \item Nếu $f$ và $g$ là các hàm số, thì $f \circ g = g \circ f$.

    \item Nếu $f$ là hàm đơn ánh, thì $f^{-1}(x) = \dfrac{1}{f(x)}$.

    \item Ta luôn luôn có thể thực hiện phép chia cho $e^x$.

    \item Nếu $0 < a < b$, thì $\ln a < \ln b$.

    \item Nếu $x > 0$, thì $(\ln x)^6 = 6\ln x$.

    \item Nếu $x > 0$ và $a > 1$, thì $\dfrac{\ln x}{\ln a} = \ln \dfrac{x}{a}$.

    \item $\tan^{-1}(-1) = 3\pi/4$.

    \item $\tan^{-1} x = \dfrac{\sin^{-1} x}{\cos^{-1} x}$.

    \item Nếu $x$ là một số thực bất kỳ, thì $\sqrt{x^2} = x$.
\end{enumerate}
\end{multicols}